\subsubsection{A 类：先交代机制，再讨论数值}

A 类行文从研究对象、坐标、尺度和守恒关系进入。高信息密度表达包括“以……为
研究对象，取……方向为正”“由质量/能量/动量守恒可得”“在边界……处满足……”
“将区域离散为……，当网格由……加密至……时，目标量变化……”。这里的关键
不是多写推导，而是使假设、方程、初边值和数值格式彼此闭合。结果解释要回到
物理量、数量级和机制，不把一条平滑曲线等同于模型正确。

需要避开的写法是先报“使用微分方程模型”，随后直接给软件输出；或者给出方程
却没有边界、单位、稳定条件和网格收敛。图形以结构示意、网格、状态曲线和残差/
收敛图为主，颜色只服务于变量区分。

